\documentclass[a4paper, 10pt]{article}

\usepackage[slovak]{babel}

% ------------------------------

\usepackage[utf8]{inputenc}
\usepackage[T1]{fontenc}

\usepackage[left=4cm,
			right=4cm,
			top=2.6cm,
			bottom=2.8cm,
			footskip=10mm,
			twoside,
			marginparwidth=3.0cm,
			%showframe,
			]{geometry}

\usepackage{graphicx}
\usepackage[dvipsnames]{xcolor}
% https://en.wikibooks.org/wiki/LaTeX/Colors

% ------------------------------

\usepackage{lmodern}

% \usepackage[tt={oldstyle=false,proportional=true,monowidth}]{cfr-lm}
\usepackage[rm={oldstyle=false}, sf={oldstyle=false}, tt={oldstyle=false,proportional=true,monowidth}]{cfr-lm}
% https://mirror.szerverem.hu/ctan/fonts/cfr-lm/doc/cfr-lm.pdf

% ------------------------------

\usepackage{amsmath}
\usepackage{amssymb}
\usepackage{amsthm}

\usepackage{booktabs}
\usepackage{multirow}
\usepackage{array}
\usepackage{dcolumn}

\usepackage{longtable}

% \usepackage{natbib}



\numberwithin{equation}{section}

% ------------------------------

\hyphenpenalty=6000
\tolerance=1000

\def\naT{\mathsf{T}}

% ------------------------------

\makeatletter

	% \def\@seccntformat#1{\protect\makebox[0pt][r]{\csname the#1\endcsname\hspace{4mm}}}

	\def\cleardoublepage{\clearpage\if@twoside \ifodd\c@page\else
	\hbox{}
	\vspace*{\fill}
	\begin{center}
	\phantom{}
	\end{center}
	\vspace{\fill}
	\thispagestyle{empty}
	\newpage
	\if@twocolumn\hbox{}\newpage\fi\fi\fi}

	\newcommand\figcaption{\def\@captype{figure}\caption}
	\newcommand\tabcaption{\def\@captype{table}\caption}

\makeatother

% ------------------------------

\usepackage{fancyhdr}
\fancypagestyle{plain}{%
\fancyhf{} % clear all header and footer fields
% \fancyfoot[C]{\sffamily {\bfseries \thepage}\ | {\scriptsize\oznacenieCasti}}
% \fancyfoot[C]{\sffamily {\bfseries \thepage}{\color{Gray}\scriptsize$\,$z$\,$\pageref{LastPage}}\ | 
\includegraphics[height=5pt]{./COMMONFILES/KUT_logo_v0.1.pdf}{\scriptsize\KUTporadoveCislo}}
\fancyfoot[C]{\sffamily {\bfseries \thepage}}
\renewcommand{\headrulewidth}{0pt}
\renewcommand{\footrulewidth}{0pt}}
\pagestyle{plain}

% ------------------------------

\usepackage{titlesec}
\titleformat{\paragraph}[hang]{\sffamily  \bfseries}{}{0pt}{}
\titlespacing*{\paragraph}{0mm}{3mm}{1mm}
\titlespacing*{\subparagraph}{0mm}{3mm}{1mm}

\titleformat*{\section}{\sffamily\Large\bfseries}
\titleformat*{\subsection}{\sffamily\large\bfseries}
\titleformat*{\subsubsection}{\sffamily\normalsize\bfseries}


% ------------------------------

\PassOptionsToPackage{hyphens}{url}
\usepackage[pdfauthor={},
			pdftitle={},
			pdfsubject={},
			pdfkeywords={},
			% hidelinks,
			colorlinks=false,
			breaklinks,
			]{hyperref}


% ------------------------------

\graphicspath{%
{./COMMONFILES/}%
{../SVG/}%
{../PY/fig/}%
{../PY/jupynotex/fig/}%
{../ML/fig/}%
{../FOTO/}%
{../SCH/DRAW}%
{../SCH/SIMULINK}%
{../SCH/TIA}%
}

% ------------------------------

\usepackage{enumitem}

\usepackage{lettrine}

% ------------------------------

\usepackage{setspace}

% ------------------------------

\usepackage{lastpage}

\usepackage{microtype}

% ------------------------------

% \usepackage[backend=biber,
%             style=numeric,
%             sorting=none,
%             ]{biblatex}
% \DeclareSourcemap{
%     \maps[datatype=bibtex]{
%         \map{
%         \step[fieldset=note, null]
%         }
%         \map{
%         \step[fieldset=file, null]
%         }        
%         % \map{
%         % \step[fieldset=url, null]        
%         % }
%         \map{
%         \step[fieldset=eprint, null]
%         }
%     }
% }
% 
% \addbibresource{./COMMONFILES/biblist.bib}



% ------------------------------

\usepackage{caption}

\usepackage{listings}




\renewcommand{\lstlistingname}{Výpis kódu}
\renewcommand{\lstlistlistingname}{Výpisy kódu}


%New colors defined below
\definecolor{codegreen}{rgb}{0,0.6,0}
\definecolor{codegray}{rgb}{0.5,0.5,0.5}
\definecolor{codepurple}{rgb}{0.58,0,0.82}
\definecolor{backcolour}{rgb}{0.95,0.95,0.95}

%Code listing style named "mystyle"
\lstdefinestyle{mystyle}{
  backgroundcolor=\color{backcolour},
  commentstyle=\fontfamily{lmtt}\fontsize{8.5pt}{8.75pt}\selectfont\color{codegreen},
  keywordstyle=\fontfamily{lmtt}\fontsize{8.5pt}{8.75pt}\selectfont\bfseries\color{Blue},
  stringstyle=\fontfamily{lmtt}\fontsize{8.5pt}{8.75pt}\selectfont\color{codepurple},
  basicstyle=\fontfamily{lmtt}\fontsize{8.5pt}{8.75pt}\selectfont,
  breakatwhitespace=false,
  breaklines=true,
  captionpos=t,
  keepspaces=true,
  numbers=left,
  numbersep=4mm,
  numberstyle=\fontfamily{lmtt}\fontsize{8.5pt}{8.75pt}\selectfont\color{lightgray},
  showspaces=false,
  showstringspaces=false,
  showtabs=false,
  tabsize=2,
  % xleftmargin=10pt,
  framesep=10pt,
  language=Python,
  escapechar=|,
}


\lstset{
    inputencoding=utf8,
    extendedchars=true,
    literate=%
    {á}{{\'a}}1
    {č}{{\v{c}}}1
    {ď}{{\v{d}}}1
    {é}{{\'e}}1
    {ě}{{\v{e}}}1
    {í}{{\'i}}1
    {ň}{{\v{n}}}1
    {ó}{{\'o}}1
    {ř}{{\v{r}}}1
    {š}{{\v{s}}}1
    {ť}{{\v{t}}}1
    {ú}{{\'u}}1
    {ů}{{\r{u}}}1
    {ý}{{\'y}}1
    {ž}{{\v{z}}}1
    {Á}{{\'A}}1
    {Č}{{\v{C}}}1
    {Ď}{{\v{D}}}1
    {É}{{\'E}}1
    {Ě}{{\v{E}}}1
    {Í}{{\'I}}1
    {Ň}{{\v{N}}}1
    {Ó}{{\'O}}1
    {Ř}{{\v{R}}}1
    {Š}{{\v{S}}}1
    {Ť}{{\v{T}}}1
    {Ú}{{\'U}}1
    {Ů}{{\r{U}}}1
    {Ý}{{\'Y}}1
    {Ž}{{\v{Z}}}1
    {ô}{{\^{o}}}1
}





\DeclareCaptionFormat{odsadene}{\protect\makebox[0pt][r]{#1#2\hspace{4mm}}#3\par}
\DeclareCaptionLabelSeparator{lendvojbodka}{:}
\DeclareCaptionFont{lightgray}{\fontfamily{lmtt}\fontsize{8.5pt}{8.75pt}\selectfont\color{lightgray}}

\captionsetup[lstlisting]{format=odsadene, labelsep=lendvojbodka, justification=raggedright, singlelinecheck=false, labelfont={sf, lightgray},}


% --- Colorful MATLAB-style listings ----------------------------------------
\lstset{
    language=Matlab,
    backgroundcolor=\color[RGB]{245,245,245}, % light gray background
    % basicstyle=\fontfamily{cmtt}\selectfont\footnotesize,
    basicstyle=\ttfamily\footnotesize,
    keywordstyle=\color[RGB]{0,0,255}\bfseries,
    stringstyle=\color[RGB]{163,21,21}, % MATLAB string color
    commentstyle=\color[RGB]{0,128,0}\itshape,
    showstringspaces=false,
    numbers=left,
    numberstyle=\scriptsize\color{gray!80},
    stepnumber=1,
    numbersep=10pt,
    frame=single,
    rulecolor=\color{gray!50},
    breaklines=true,
    breakatwhitespace=true,
    tabsize=4,
    captionpos=b,
    emph={for,end,if,else,elseif,while,function,return,case,otherwise,try,catch},
    emphstyle={\color[RGB]{0,0,255}\bfseries}, % MATLAB keywords
    columns=fullflexible,
    keepspaces=true,
    upquote=true
}

% ------------------------------

\lstdefinestyle{pseudocode}{
    backgroundcolor=\color{white},
    basicstyle=\ttfamily\small,
    keywordstyle=\bfseries,
    commentstyle=\itshape\color{gray},
    stringstyle=\color{black},
    numbers=none,
    frame=single,
    rulecolor=\color{black!20},
    breaklines=true,
    tabsize=2,
    captionpos=b,
    columns=fullflexible,
    keepspaces=true,
    showstringspaces=false,
    morekeywords={INPUT,OUTPUT,IF,THEN,ELSE,WHILE,FOR,RETURN,END},
    morecomment=[l]{\%}
}

% Cool cpp style------------------------------




\definecolor{codebg}{RGB}{30,32,34}
\definecolor{codeframe}{RGB}{70,70,70}
\definecolor{codekw}{RGB}{86,156,214}
\definecolor{codestr}{RGB}{214,157,133}
\definecolor{codecom}{RGB}{106,153,85}
\definecolor{codenum}{RGB}{120,120,120}
\definecolor{codeid}{RGB}{220,220,220}
\definecolor{codeemph}{RGB}{78,201,176}

\lstdefinestyle{cppstyle}{
    language=C++,
    backgroundcolor=\color{codebg},
    basicstyle=\ttfamily\footnotesize\color{codeid},
    keywordstyle=\color{codekw}\bfseries,
    stringstyle=\color{codestr},
    commentstyle=\color{codecom}\itshape,
    identifierstyle=\color{codeid},
    emphstyle=\color{codeemph}\bfseries,
    numbers=left,
    numberstyle=\scriptsize\color{codenum},
    stepnumber=1,
    numbersep=8pt,
    frame=none,
    breaklines=true,
    breakatwhitespace=false,
    tabsize=4,
    captionpos=b,
    columns=fixed,
    keepspaces=true,
    showstringspaces=false,
    xleftmargin=0pt,
    xrightmargin=0pt,
    aboveskip=4pt,
    belowskip=4pt
}


\begin{document}


% -------------------------------------------------------------------
% Titulný list - Diplomová práca STU FEI
% -------------------------------------------------------------------


\thispagestyle{empty}

\begin{center}
    {\bfseries\Large SLOVENSKÁ TECHNICKÁ UNIVERZITA V BRATISLAVE}\\[3pt]
    {\bfseries\large Fakulta elektrotechniky a informatiky}\\[5mm]

    {\bfseries Evidenčné číslo: FEI-104376-117218}\\[80mm]

    {\bfseries\LARGE Návrh implementácie vybraných algoritmov riadenia}\\[20pt]
    {\large\ Diplomová práca}\\[55mm]
\end{center}

\begin{flushleft}
    {\bfseries Študijný program:}\hspace{34mm} Robotika a kybernetika\\[4pt]
    {\bfseries Študijný odbor:}\hspace{39.5mm} kybernetika\\[4pt]
    {\bfseries Školiace pracovisko:}\hspace{31mm} Ústav robotiky a kybernetiky\\[4pt]
    {\bfseries Vedúci záverečnej práce/školiteľ:}\hspace{4.5mm} Ing. Marián Tárnik, PhD.\\[25mm]
\end{flushleft}

\begin{flushleft}
    {\bfseries Bratislava 2026}\hspace{60mm}Bc. Andrii Stoliar
\end{flushleft}

\clearpage
\setcounter{page}{1}

% -------------------------------------------------------------------
% Strana s obrazkom zadania DP
% -------------------------------------------------------------------

\clearpage
\thispagestyle{empty}

\noindent
\hspace*{-\dimexpr\oddsidemargin+1in\relax}%
\makebox[\paperwidth][l]{%
  \includegraphics[
    width=\paperwidth,        
    keepaspectratio
  ]{zadanie.png}%
}

\clearpage

% -------------------------------------------------------------------
% Poďakovanie (3. strana)
% -------------------------------------------------------------------

\clearpage

\thispagestyle{empty}
\begin{center}
    {\bfseries\LARGE Poďakovanie}\\[170mm]
\end{center}

\noindent
Na tomto mieste by som rád poďakoval svojmu školiteľovi Ing. Mariánovi Tárnikovi, PhD.,
za odborné vedenie, cenné rady a trpezlivosť počas vypracovania tejto diplomovej práce.\\
Zároveň ďakujem celému Ústavu robotiky a kybernetiky za odovzdanie cenných vedomostí a skúseností,
ktoré mi pomohli rozvíjať sa v oblasti riadenia a automatizácie.\\
Veľká vďaka patrí aj mojim priateľom, starým rodičom, mame, otcovi, sestre a mojej dáme srdca
za ich obrovskú podporu, povzbudenie a pochopenie počas celého štúdia.\\[5mm]

\begin{flushright}
    \textit{Andrii Stoliar}\\
    Bratislava, 2026
\end{flushright}

\clearpage
% -------------------------------------------------------------------
% Obsah (6. strana)
% -------------------------------------------------------------------

\clearpage
\thispagestyle{empty}

\tableofcontents

\clearpage

% -------------------------------------------------------------------
% Uvod
% -------------------------------------------------------------------

\clearpage
\pagenumbering{arabic}
\setcounter{page}{1}
\pagestyle{fancy}

\fancyhf{}
\fancyfoot[C]{\thepage\ z \pageref{LastPage}}
\renewcommand{\headrulewidth}{0pt}

\section*{Úvod}
\addcontentsline{toc}{section}{Úvod}
Riadenie procesov je kľúčovou súčasťou moderných priemyselných systémov,
ktoré zabezpečujú efektívnu a spoľahlivú prevádzku rôznych technológií.
   
\clearpage

% -------------------------------------------------------------------
% Section 1 - Tepelny system
% -------------------------------------------------------------------

\clearpage

\section{Tepelný systém}

V tejto kapitole sa budeme zaoberať laboratórnym modelom, ktorý bol poskytnutý pre potreby tejto diplomovej práce – tepelným systémom TS05.  
Cieľom kapitoly je podrobne opísať túto experimentálnu zostavu, vykonať jej identifikáciu a navrhnúť viaceré riešenia riadenia, a to ako klasické, tak aj pokročilejšie metódy.

Súčasťou práce je tiež implementácia vybraných riadiacich algoritmov na rôznych hardvérových platformách.  
Táto časť má umožniť porovnanie prístupov k implementácii riadenia v simulačnom prostredí \textit{MATLAB/Simulink} a na vybranom mikrokontroléri, čím sa vytvorí priestor pre analýzu rozdielov v správaní systému a v samotnej realizácii algoritmov.

% -------------------------------------------------------------------
% Subsection 1.1 - Popis tepelneho systemu
% -------------------------------------------------------------------

\subsection{Popis systému TS05}

Laboratórny model TS05 predstavuje \textit{MIMO} systém (Multiple Input – Multiple Output) so dvomi vstupmi a dvomi výstupmi, ktorého zapojenie je uvedené v tabuľke~\ref{tab:io_ts05}.  
Vstupné veličiny tvoria napätia riadiace \textit{ventilátor} a \textit{vykurovaciu špirálu} (v rozsahu 0–10~V), zatiaľ čo výstupy sú merané pomocou \textit{Senzora~1} a \textit{Senzora~2} (tiež v rozsahu 0–10~V).

\begin{table}[h!]
    \centering
    \caption{Vstupy a výstupy systému TS05}
    \label{tab:io_ts05}
    \begin{tabular}{|c|c|c|}
        \hline
        \textbf{Typ signálu} & \textbf{Označenie} & \textbf{Rozsah} \\ \hline
        Vstup & Ventilátor & 0–10~V \\ \hline
        Vstup & Výkon špirály & 0–10~V \\ \hline
        Výstup & Senzor~1 (teplota 1) & 0–10~V \\ \hline
        Výstup & Senzor~2 (teplota 2) & 0–10~V \\ \hline
    \end{tabular}
\end{table}

Vizuálny vzhľad systému je znázornený na obrázku~\ref{fig:ts05_real},  
kde sú zreteľne vyznačené všetky vstupy a výstupy zariadenia.

\begin{center}
\vbox{%
\makebox[\textwidth][c]{%
\includegraphics[width=1.1\textwidth]{TS_sch3_foto.pdf}
}
\figcaption{Laboratórny model TS05 – reálna zostava}
\label{fig:ts05_real}
}%vbox
\end{center}

Systém komunikuje s počítačom prostredníctvom meracej a komunikačnej karty \textit{Humusoft},  
ktorá je určená na prepojenie reálnych systémov so simulačným prostredím \textit{MATLAB/Simulink}.  
Táto karta umožňuje obojsmerný prenos analógových signálov – prijíma merané veličiny zo systému a zároveň vysiela riadiace signály na jeho vstupy.

Samotná komunikácia a spracovanie signálov prebieha v prostredí \textit{Simulink},  
v rámci vopred pripravenej schémy (obr.~\ref{fig:simulink_schema}),  
ktorá bola vytvorená ešte pred začiatkom tohto akademického roka ako súčasť projektu KUT – Krátke výučbové materiály.  
Táto schéma slúžila aj ako podklad pre predmet \textit{Modelovanie a riadenie systémov}.

\begin{center}
\vbox{%
\makebox[\textwidth][c]{%
\includegraphics[width=1.3\textwidth, trim=100 200 100 200, clip]{TS_simulink_scheme.pdf}
}
\figcaption{Simulinková schéma pre komunikáciu so systémom TS05}
\label{fig:simulink_schema}
}%vbox
\end{center}

% -------------------------------------------------------------------
% Subsection 1.2 - Identifikacia tepelneho systemu
% -------------------------------------------------------------------

\subsection{Identifikácia systému TS05}

Keďže ide o \textin{MIMO systém} (Multiple Input – Multiple Output), jeho identifikácia v jednom kroku nie je realizovateľná.  
Z tohto dôvodu bolo rozhodnuté vykonať identifikáciu systému osobitne pre tri pracovné body ventilátora – 3~V, 6~V a 9~V – a pre oba senzory samostatne.  
Takto získame komplexnejší pohľad na správanie systému v rôznych prevádzkových režimoch.

% -------------------------------------------------------------------
% Subsubsection 1.2.1 - Merabie prechodovych charakteristik
% -------------------------------------------------------------------

\subsubsection{Meranie prechodových charakteristík}

Ako bolo spomenuté, prvým krokom bola identifikácia systému prostredníctvom merania jeho statických a prechodových charakteristík.  
Na tento účel boli vykonané tri experimenty, zodpovedajúce trom pracovným bodom ventilátora (3~V, 6~V a 9~V).  
Počas každého experimentu bola hodnota napätia ventilátora udržiavaná konštantná, zatiaľ čo napätie vykurovacej špirály sa lineárne zvyšovalo z 0~V na 10~V počas celého trvania merania, keďže špirála má výrazne dynamickejší vplyv na teplotu ako ventilátor.  
Priebeh signálu špirály počas experimentov je znázornený na obrázku~\ref{fig:spiral_signal}.

\begin{center}
\vbox{%
\makebox[\textwidth][c]{%
\includegraphics[width=1.5\textwidth]{ts_prevodova_spirala.png}
}
\figcaption{Priebeh napätia na špirále počas všetkých troch experimentov}
\label{fig:spiral_signal}
}%vbox
\end{center}

Výsledkom meraní sú šesť grafov, ktoré znázorňujú odozvu systému na zmenu vstupných veličín.  
Každý graf zobrazuje výstupné signály zo Senzora~1 a Senzora~2 ako funkciu času pre daný pracovný bod ventilátora.

\begin{center}
\vbox{%
\makebox[\textwidth][c]{%
\includegraphics[width=1.5\textwidth]{ts_prevodova_sensors3.png}
}
\figcaption{Odozva systému TS05 pre pracovný bod ventilátora 3~V}
\label{fig:transient_3V}
}%vbox
\end{center}

\begin{center}
\vbox{%
\makebox[\textwidth][c]{%
\includegraphics[width=1.5\textwidth]{ts_prevodova_sensors6.png}
}
\figcaption{Odozva systému TS05 pre pracovný bod ventilátora 6~V}
\label{fig:transient_6V}
}%vbox
\end{center}

\begin{center}
\vbox{%
\makebox[\textwidth][c]{%
\includegraphics[width=1.5\textwidth]{ts_prevodova_sensors9.png}
}
\figcaption{Odozva systému TS05 pre pracovný bod ventilátora 9~V}
\label{fig:transient_9V}
}%vbox
\end{center}

Na účely kvalitnejšieho opisu systému z hľadiska statických charakteristík boli z rovnakých meraní zostavené aj prevodové charakteristiky, ktoré opisujú závislosť výstupných veličín od vstupného napätia špirály pre jednotlivé pracovné body ventilátora.

\begin{center}
\vbox{%
\makebox[\textwidth][c]{%
\includegraphics[width=1.5\textwidth]{ts_static_sensors3.png}
}
\figcaption{Prevodová charakteristika systému TS05 pre pracovný bod ventilátora 3~V}
\label{fig:static_3V}
}%vbox
\end{center}

\begin{center}
\vbox{%
\makebox[\textwidth][c]{%
\includegraphics[width=1.5\textwidth]{ts_static_sensors6.png}
}
\figcaption{Prevodová charakteristika systému TS05 pre pracovný bod ventilátora 6~V}
\label{fig:static_6V}
}%vbox
\end{center}

\begin{center}
\vbox{%
\makebox[\textwidth][c]{%
\includegraphics[width=1.5\textwidth]{ts_static_sensors9.png}
}
\figcaption{Prevodová charakteristika systému TS05 pre pracovný bod ventilátora 9~V}
\label{fig:static_9V}
}%vbox
\end{center}

% -------------------------------------------------------------------
% Subsubsection 1.2.2 - Identifikacia pomocou metody najmensich stvorcov
% -------------------------------------------------------------------

\subsubsection{Identifikácia pomocou metódy najmenších štvorcov}

V tejto časti bola vykonaná identifikácia všetkých šiestich nameraných odoziev systému TS05.  
Postup identifikácie bol navrhnutý nasledovne: z každej prechodovej charakteristiky bolo vybraných sedem počiatočných skokov, pretože po tomto bode systém vo väčšine prípadov dosahuje stav nasýtenia, čo je možné pozorovať aj na predošlých grafoch odozvy systému.  
Pre každé meranie bol z týchto siedmich odoziev vypočítaný ich priemerný priebeh, čím sa získal reprezentatívny „priemerný skok“ charakterizujúci správanie systému v danej prevádzkovej oblasti.

Na základe tohto priemerného skoku bola následne uskutočnená identifikácia systému pomocou metódy najmenších štvorcov~\cite{fikar}.  
Táto metóda umožňuje odhad parametrov diskrétneho modelu systému na základe meraných vstupov a výstupov.  
Pre každý z experimentov bola vytvorená regresná matica $\Phi(k,:)$ a vektor výstupov $Y$, pričom výsledný odhad parametrov $\hat{\theta}$ bol určený podľa rovníc:

\begin{equation}
    \Phi(k,:) = [-y(k-1), -y(k-2), u(k-1), u(k-2)]
    \label{eq:phi_matrix}
\end{equation}

\begin{equation}
    \hat{\theta} = (\Phi^\mathrm{T}\Phi)^{-1}\Phi^\mathrm{T}Y
    \label{eq:theta_estimation}
\end{equation}

V prostredí \textit{MATLAB} bola implementovaná funkcia, ktorá realizuje uvedené výpočty 
a vytvára diskrétny model systému $G_d$, ktorý je následne prevedený na spojitú 
prenosovú funkciu $G$ metódou \textit{Tustinovej aproximácie}.  
Implementácia tejto funkcie je uvedená nižšie:

\begin{lstlisting}[language=Matlab, 
caption={MATLAB implementácia identifikácie metódou najmenších štvorcov}, 
label={lst:ls_identify}]
function G = ls_identify_basic(y, u, Ts)

    y = y(:);
    u = u(:);
    N = length(y);
    na = 2; nb = 2;
    
    Phi = zeros(N-2, na+nb);
    for k = 3:N
        Phi(k-2,:) = [-y(k-1), -y(k-2), u(k-1), u(k-2)];
    end
    Y = y(3:N);
    
    theta = pinv(Phi) * Y;
    
    a1 = theta(1);
    a2 = theta(2);
    b1 = theta(3);
    b2 = theta(4);
    
    num = [b1 b2];
    den = [1 a1 a2];
    Gd = tf(num, den, Ts);
    
    G = d2c(Gd, 'tustin');
end
\end{lstlisting}

Následne bol spustený identifikačný skript pre všetkých šesť meraných priebehov.  
Porovnanie medzi reálnymi priemernými skokovými odozvami a odozvami identifikovaných modelov 
je zobrazené na obrázku~\ref{fig:ts_idented_steps}.  

Z porovnania možno pozorovať, že počas 100-sekundového merania senzor~2, 
ktorý je umiestnený ďalej od zdroja tepla, nedosiahne úplne ustálenú hodnotu.  
Tento jav je však možné považovať za zanedbateľný, keďže jeho vplyv na kvalitu 
následného návrhu riadenia je minimálny — model dostatočne zachytáva dynamické správanie systému 
v oblasti, ktorá je relevantná pre reguláciu.

Pre účely ďalšej analýzy a návrhu regulátora bolo rozhodnuté aproximovať správanie systému 
modelom so dvoma pólmi a dvoma nulami.  
Takýto zjednodušený model je výhodný nielen z pohľadu výpočtovej efektívnosti, 
ale aj z hľadiska laditeľnosti parametrov riadiaceho algoritmu.  
Zníženie poradia systému zaisťuje lepšiu interpretovateľnosť a stabilnejšie správanie modelu 
pri simuláciách aj pri následnej implementácii na reálnom hardvéri.

\begin{center}
\vbox{%
\makebox[\textwidth][c]{%
\includegraphics[width=1.25\textwidth]{ts_idented_steps.png}
}
\figcaption{Porovnanie reálnych a identifikovaných skokových odoziev systému TS05}
\label{fig:ts_idented_steps}
}%vbox
\end{center}

\begin{table}[h!]
\centering
\caption{Prenosové funkcie identifikované pre jednotlivé pracovné body systému TS05}
\label{tab:identified_models}
\renewcommand{\arraystretch}{1.3}
\begin{tabular}{|c|c|l|}
\hline
\textbf{Pracovný bod [V]} & \textbf{Senzor} & \textbf{Prenosová funkcia $G(s)$} \\ \hline
3 & Near &
$\displaystyle G_{3N}(s) = \frac{1.385\times10^{-15}s^2 - 0.1973s + 3.947}{s^2 + 67.03s + 5.015}$ \\ \hline
3 & Far &
$\displaystyle G_{3F}(s) = \frac{-2.898\times10^{-18}s^2 - 0.1044s + 2.088}{s^2 + 56s + 3.534}$ \\ \hline
6 & Near &
$\displaystyle G_{6N}(s) = \frac{-5.539\times10^{-17}s^2 - 0.2101s + 4.201}{s^2 + 69.59s + 6.914}$ \\ \hline
6 & Far &
$\displaystyle G_{6F}(s) = \frac{1.165\times10^{-15}s^2 - 0.1216s + 2.431}{s^2 + 59.73s + 5.338}$ \\ \hline
9 & Near &
$\displaystyle G_{9N}(s) = \frac{-2.405\times10^{-17}s^2 - 0.1969s + 3.938}{s^2 + 64.82s + 6.505}$ \\ \hline
9 & Far &
$\displaystyle G_{9F}(s) = \frac{-4.79\times10^{-16}s^2 - 0.113s + 2.26}{s^2 + 58.73s + 5}$ \\ \hline
\end{tabular}
\end{table}

\vspace{4mm}

Na základe identifikácie sme získali šesť prenosových funkcií, 
ktoré reprezentujú správanie systému pri rôznych pracovných bodoch ventilátora 
(3~V, 6~V a~9~V) a~pre obe meracie pozície (\textit{near} – bližší senzor, 
\textit{far} – vzdialenejší senzor).  
Tieto modely môžu slúžiť ako podklad pre návrh a~ladenie parametrov \textit{PID} regulátorov.

Tieto prenosové funkcie poskytujú základ pre návrh regulácie, 
umožňujú odhadnúť dynamiku systému v~rôznych prevádzkových režimoch 
a~zároveň slúžia na porovnanie vplyvu vzdialenosti senzora na odozvu systému.

% -------------------------------------------------------------------
% Zoznam literatúry 
% -------------------------------------------------------------------

\clearpage

\addcontentsline{toc}{section}{Zoznam literatúry}

\begin{thebibliography}{9}

\bibitem{fikar}
M. Fikar, J. Mikleš, \textit{Identifikácia systémov}.  
Slovenská technická univerzita v Bratislave, 1998.  
ISBN 80-227-1177-2.  
Dostupné online: \url{http://TODO.htm}

\end{thebibliography}


\makeatletter
\@ifundefined{r@LastPage}{\newlabel{LastPage}{{}{\thepage}}}{}
\makeatother

\clearpage


\end{document}
